% Submission Site: https://mc.manuscriptcentral.com/taco
% Guidelines: https://dl.acm.org/journal/taco/author-guidelines

% Page Limit:
%%   Original submission: 20 pages using the 7x10 ACM Submission Format. The 20 pages includes text, figures, but no references. This is a strict limit for the initial submission.
%%   Revised versions of the manuscript: 25 pages, including references.

% Other requirements:
%%  ORCID Requirements
%%      ACM requires that all accepted journal authors register and provide ACM with valid ORCIDs prior to paper publication.

\PassOptionsToPackage{table,dvipsnames}{xcolor}
% \documentclass[manuscript,screen,review,numbers,sort&compress,hyphens]{acmart}
\documentclass[manuscript,screen,review]{acmart}
% \documentclass[sigconf,screen,review]{acmart}
\citestyle{acmnumeric}

\input{local}
\input{mycommand}

% Keep the revised manuscript within the 25-page limit without changing the
% body text or the revision-marked layout.  This preserves acmart's public
% section commands, fonts, and numbering and only substitutes slightly tighter
% spacing when its internal heading formatter receives a section or subsection.
\makeatletter
\let\TilingInfer@startsection\@startsection
\def\@startsection#1#2#3#4#5#6{%
  \def\TilingInfer@sectionname{#1}%
  \def\TilingInfer@section{section}%
  \def\TilingInfer@subsection{subsection}%
  \def\TilingInfer@beforeskip{#4}%
  \def\TilingInfer@afterskip{#5}%
  \ifx\TilingInfer@sectionname\TilingInfer@section
    \def\TilingInfer@beforeskip{-.68\baselineskip
      \@plus -2\p@ \@minus -.2\p@}%
    \def\TilingInfer@afterskip{.20\baselineskip}%
  \else
    \ifx\TilingInfer@sectionname\TilingInfer@subsection
      \def\TilingInfer@beforeskip{-.68\baselineskip
        \@plus -2\p@ \@minus -.2\p@}%
      \def\TilingInfer@afterskip{.20\baselineskip}%
    \fi
  \fi
  \TilingInfer@startsection{#1}{#2}{#3}%
    {\TilingInfer@beforeskip}{\TilingInfer@afterskip}{#6}}
\makeatother

%%
%% \BibTeX command to typeset BibTeX logo in the docs
% \AtBeginDocument{%
%   \providecommand\BibTeX{{%
%     Bib\TeX}}}

%% Rights management information.  This information is sent to you
%% when you complete the rights form.  These commands have SAMPLE
%% values in them; it is your responsibility as an author to replace
%% the commands and values with those provided to you when you
%% complete the rights form.
% \setcopyright{acmlicensed}
% \copyrightyear{2018}
% \acmYear{2018}
% \acmDOI{XXXXXXX.XXXXXXX}
% %% These commands are for a PROCEEDINGS abstract or paper.
% \acmConference[Conference acronym 'XX]{Make sure to enter the correct
%   conference title from your rights confirmation email}{June 03--05,
%   2018}{Woodstock, NY}
% %%
% %%  Uncomment \acmBooktitle if the title of the proceedings is different
% %%  from ``Proceedings of ...''!
% %%
% %%\acmBooktitle{Woodstock '18: ACM Symposium on Neural Gaze Detection,
% %%  June 03--05, 2018, Woodstock, NY}
% \acmISBN{978-1-4503-XXXX-X/2018/06}


%%
%% Submission ID.
%% Use this when submitting an article to a sponsored event. You'll
%% receive a unique submission ID from the organizers
%% of the event, and this ID should be used as the parameter to this command.
%%\acmSubmissionID{123-A56-BU3}

%%
%% For managing citations, it is recommended to use bibliography
%% files in BibTeX format.
%%
%% You can then either use BibTeX with the ACM-Reference-Format style,
%% or BibLaTeX with the acmnumeric or acmauthoryear sytles, that include
%% support for advanced citation of software artefact from the
%% biblatex-software package, also separately available on CTAN.
%%
%% Look at the sample-*-biblatex.tex files for templates showcasing
%% the biblatex styles.
%%

%%
%% The majority of ACM publications use numbered citations and
%% references.  The command \citestyle{authoryear} switches to the
%% "author year" style.
%%
%% If you are preparing content for an event
%% sponsored by ACM SIGGRAPH, you must use the "author year" style of
%% citations and references.
%% Uncommenting
%% the next command will enable that style.
%%\citestyle{acmauthoryear}


%%
%% end of the preamble, start of the body of the document source.
\begin{document}

%%
%% The "title" command has an optional parameter,
%% allowing the author to define a "short title" to be used in page headers.
\title{Optimizing Dynamic-Shape Operator Libraries via Automatic Model-Specific Specialization on AI Accelerators}

%%
%% The "author" command and its associated commands are used to define
%% the authors and their affiliations.
%% Of note is the shared affiliation of the first two authors, and the
%% "authornote" and "authornotemark" commands
%% used to denote shared contribution to the research.
% \author{Shuo Liu}
% \authornote{Both authors contributed equally to this research.}
% \email{zkdliushuo@mail.ustc.edu.cn}
% \orcid{1234-5678-9012}
% \affiliation{%
%   \institution{University of Science and Technology of China}
%   \city{Hefei}
%   \state{Anhui}
%   \country{CN}
% }

% \author{Lars Th{\o}rv{\"a}ld}
% \affiliation{%
%   \institution{The Th{\o}rv{\"a}ld Group}
%   \city{Hekla}
%   \country{Iceland}}
% \email{larst@affiliation.org}

%%
%% By default, the full list of authors will be used in the page
%% headers. Often, this list is too long, and will overlap
%% other information printed in the page headers. This command allows
%% the author to define a more concise list
%% of authors' names for this purpose.
% \renewcommand{\shortauthors}{Shuo et al.}

\input{sections/abstract.tex}

%%
%% The code below is generated by the tool at http://dl.acm.org/ccs.cfm.
%% Please copy and paste the code instead of the example below.
\begin{CCSXML}
	<ccs2012>
	<concept>
		<concept_id>10011007.10011006.10011041.10011045</concept_id>
		<concept_desc>Software and its engineering~Dynamic compilers</concept_desc>
		<concept_significance>500</concept_significance>
		</concept>
	<concept>
		<concept_id>10010520.10010521.10010528.10010534</concept_id>
		<concept_desc>Computer systems organization~Single instruction, multiple data</concept_desc>
		<concept_significance>300</concept_significance>
		</concept>
	<concept>
		<concept_id>10010147.10010257.10010293.10010294</concept_id>
		<concept_desc>Computing methodologies~Neural networks</concept_desc>
		<concept_significance>300</concept_significance>
		</concept>
	<concept>
		<concept_id>10002944.10011123.10011674</concept_id>
		<concept_desc>General and reference~Performance</concept_desc>
		<concept_significance>300</concept_significance>
		</concept>
	</ccs2012>
\end{CCSXML}

\ccsdesc[500]{Software and its engineering~Dynamic compilers}
\ccsdesc[300]{Computer systems organization~Single instruction, multiple data}
\ccsdesc[300]{Computing methodologies~Neural networks}
\ccsdesc[300]{General and reference~Performance}

\keywords{compilers, static analysis, deep learning, kernel specialization, constant propagation, AI accelerators, operator libraries}

% \received{20 February 2007}
% \received[revised]{12 March 2009}
% \received[accepted]{5 June 2009}

%%
%% This command processes the author and affiliation and title
%% information and builds the first part of the formatted document.
\maketitle

\input{sections/introduction}
\input{sections/background.tex}
\input{sections/design.tex}
\input{sections/evaluation.tex}
\input{sections/related.tex}
\section{Discussion}
\label{sec:discussion}

\textbf{Workload-dependent benefits.}
Our experiments show that the magnitude of \mysys's benefit depends strongly on
workload characteristics. Recommendation and decode workloads benefit most
because Attention and MatMul expose profitable invariant schedules, yielding
about 10\% average inference speedup. Schedule-fixed expert-written libraries
also benefit from removing unreachable template instances. This per-model
pruning is most attractive when a deployment serves fewer models, because each
pruned binary can discard a larger fraction of the universal library. By
contrast, HSTU and prefill remain compute-heavy, so specializing their core
Attention and MatMul kernels removes little of the total work and improves
end-to-end latency by only about 3\%. Some kernels in the current open-source
baselines also cannot be further specialized by \mysys, further limiting
end-to-end gains.

\textbf{Platform applicability.}
Our schedule-generic kernel-specialization implementation currently targets
Ascend. In the open-source libraries we surveyed for other accelerators,
expert-written kernels predominantly follow the schedule-fixed design, leaving
no comparable schedule-generic baseline to evaluate. Ascend exposes many
programmable hardware levels; achieving high performance therefore relies on
numerous runtime schedule parameters, which motivates schedule-generic
libraries and makes automatic specialization particularly relevant.

\textbf{Optimization scope and limitations.}
\mysys applies to existing expert-written operator libraries: it specializes
schedule-generic kernels or removes unreachable instances from schedule-fixed
libraries. This differs from AI compilers, which generate implementations and
search schedules; their code-generation pipelines do not optimize an existing
C++ expert library in place. 

In addition, \mysys analyzes host code only to recover constants
and launch-site liveness, then optimizes reachable device kernels; it does not
rewrite host code. Host control flow is complex, CPUs already execute it
efficiently, and profitable host specialization would require a more selective
policy than device-kernel constant propagation. NPU Graph replay suppresses
repeated launch overhead, while identical input shapes across Transformer
layers allow host-computed schedule parameters to be reused. Host schedule
computation is therefore not a bottleneck in our evaluated setting, making
device-kernel optimization the higher-value target.

\input{sections/conclusion}

%%
%% The next two lines define the bibliography style to be used, and
%% the bibliography file.
\bibliographystyle{ACM-Reference-Format}
\bibliography{opt,ref,pa}

\end{document}
\endinput
%%
%% End of file `sample-manuscript.tex'.
